% =============================================================
% homework.tex — Домашняя работа по математике
% Формат: A4 portrait
% Компилятор: XeLaTeX
%
% Структура страницы:
%
%   Домашняя работа              Класс: ___
%   Тема: Дроби и смешанные числа
%   Имя и фамилия: _________ Сдать до: 14 марта
%   ───────────────────────────────────────────
%
%   Упражнения
%
%   1.  3/4 + 5/8 =  ______    2. ...    3. ...
%
%   4.  [задание с 2 строками]
%
%   5.  [задача с полем для решения]
%       |
%       |  ________________________________
%       |  ________________________________
%       |  ________________________________
%       Ответ: _______________
%
%   ◆  Напоминание: a/b + c/b = (a+c)/b
%
%   ───────────────────────────────────────────
%   Проверил: _________________  Подпись: ______
% =============================================================

\documentclass[12pt]{article}

\usepackage{../styles/mathworksheet}

% Переключаем на геометрию домашней работы (portrait, поля 25/15/18/18)
\hwsetgeometry

% ============================================================
% ПАРАМЕТРЫ — ИЗМЕНЯЙТЕ ЗДЕСЬ
% ============================================================
\newcommand{\hwtopic}{Сложение и вычитание дробей}
\newcommand{\hwdeadline}{14 марта}
% ============================================================

\begin{document}

% Ширина текстовой колонки заданий
\settaskwidth{155mm}

% ============================================================
% ШАПКА
% ============================================================
\hwheader{}{\hwtopic}{\hwdeadline}

% ============================================================
% РАЗДЕЛ 1: ВЫЧИСЛИ
% ============================================================
\hwsection{Упражнения}

% Три вычислительных в ряд через minipage (горизонтальная раскладка)
\noindent
\begin{minipage}[t]{56mm}
  \hwcalc{1}{\tfrac{3}{4} + \tfrac{5}{8}}
\end{minipage}%
\begin{minipage}[t]{56mm}
  \hwcalc{2}{\tfrac{7}{6} - \tfrac{2}{3}}
\end{minipage}%
\begin{minipage}[t]{56mm}
  \hwcalc{3}{2\tfrac{1}{4} + 1\tfrac{3}{8}}
\end{minipage}

\vspace{1mm}

\noindent
\begin{minipage}[t]{56mm}
  \hwcalc{4}{\tfrac{5}{9} \cdot \tfrac{3}{10}}
\end{minipage}%
\begin{minipage}[t]{56mm}
  \hwcalc{5}{3\tfrac{1}{2} - 1\tfrac{5}{6}}
\end{minipage}%
\begin{minipage}[t]{56mm}
  \hwcalc{6}{\tfrac{4}{5} : \tfrac{8}{15}}
\end{minipage}

\vspace{2mm}

% Подсказка для раздела
\hwhint{Чтобы сложить дроби с разными знаменателями, приведите их к общему знаменателю.}

% ============================================================
% РАЗДЕЛ 2: ЗАДАНИЯ С ЗАПИСЬЮ РЕШЕНИЯ
% ============================================================
\hwsection{Задания}

\hwtasklines{7}{%
  Упростите выражение:\quad
  $\dfrac{5}{12} + \dfrac{3}{8} - \dfrac{1}{6}$%
}{3}

\hwtasklines{8}{%
  Найдите значение выражения при $a = \tfrac{2}{3}$ и $b = \tfrac{3}{4}$:\quad
  $\dfrac{a + b}{2}$%
}{3}

% ============================================================
% РАЗДЕЛ 3: ЗАДАЧА
% ============================================================
\hwsection{Задача}

\hwproblem{9}{%
  На приготовление обеда ушло $\tfrac{3}{4}$~ч, а на
  приготовление ужина~— на $\tfrac{1}{6}$~ч меньше.
  Сколько времени в сумме потратили на готовку?%
}{3}

% ============================================================
% НИЖНИЙ КОЛОНТИТУЛ
% ============================================================
\vfill
\hwfooter

\end{document}

% ============================================================
% СПРАВОЧНИК КОМАНД
% ============================================================
%
% Шапка:
%   \hwheader{}{тема}{дата сдачи}
%     Первый аргумент зарезервирован, оставляйте пустым —
%     класс ученик пишет сам.
%
% Разделы:
%   \hwsection{название}         — серый курсивный заголовок
%
% Задания:
%   \hwcalc{N}{выражение}        — вычислительное (= ______)
%   \hwtask{N}{текст}            — задание с воздухом под ним
%   \hwtasklines{N}{текст}{L}    — задание + L строк для записи
%   \hwproblem{N}{условие}{L}    — задача + L строк + «Ответ:»
%
% Подсказка:
%   \hwhint{текст}               — серый курсив с ромбом
%
% Горизонтальный ряд вычислений — через minipage:
%   \noindent
%   \begin{minipage}[t]{56mm}\hwcalc{1}{...}\end{minipage}%
%   \begin{minipage}[t]{56mm}\hwcalc{2}{...}\end{minipage}%
%   \begin{minipage}[t]{56mm}\hwcalc{3}{...}\end{minipage}
%   (3 × 54mm = 162mm ≈ \hwtextwidth)
%
% Нижний колонтитул:
%   \vfill\hwfooter
%
% КОМПИЛЯЦИЯ:
%   xelatex homework.tex
% ============================================================